\section{Problem 3: Three-Axis Gyroscope}
\begin{figure} [h]
\centering
\includegraphics[width=0.7\textwidth]{figures/Problem_3.png}
\caption{Three-axis gyroscope system.}
\end{figure}
We use the same Euler angles and reference frames from HW01:
\[
q_1 = \psi,
\qquad
q_2 = \theta,
\qquad
q_3 = \phi
\]
where:
\begin{itemize}
    \item $\psi$ is the precession angle,
    \item $\theta$ is the nutation angle,
    \item $\phi$ is the rotor spin angle.
\end{itemize}

From HW01, the angular velocity expressed in the rotor-fixed frame $B$ is:
\[
\bm{\omega}_{B/N}
=
\dot\theta\mathbf{b}_1
+
\dot\psi\sin\theta\mathbf{b}_2
+
(\dot\phi + \dot\psi\cos\theta)\mathbf{b}_3
\]

\subsection*{(a) Inertia Tensor About $O$}

The rotor is modeled as a uniform axisymmetric disk of radius $R$ and mass $m$.

The principal moments of inertia are:
\[
I_1 = I_2 = \frac14 mR^2,
\qquad
I_3 = \frac12 mR^2
\]

Thus the inertia tensor about the fixed point $O$ in frame $B$ is:
\[
\mathbf{I}_O
=
\begin{bmatrix}
I_1 & 0 & 0 \\
0 & I_1 & 0 \\
0 & 0 & I_3
\end{bmatrix}
\]

\subsection*{(b) Angular Momentum About $O$}

The angular momentum is:
\[
\mathbf{H}_O = \mathbf{I}_O\bm{\omega}_{B/N}
\]

Substituting the angular velocity:
\begin{align}
\mathbf{H}_O
&=
I_1\dot\theta\mathbf{b}_1
+
I_1\dot\psi\sin\theta\mathbf{b}_2
\\
&\quad
+
I_3(\dot\phi+\dot\psi\cos\theta)\mathbf{b}_3
\end{align}

Hence the components are:
\begin{align}
H_1 &= I_1\dot\theta
\\
H_2 &= I_1\dot\psi\sin\theta
\\
H_3 &= I_3(\dot\phi+\dot\psi\cos\theta)
\end{align}

\subsection*{(c) Euler's Equations}

Since point $O$ is fixed in the inertial frame:
\[
\sum \mathbf{M}_O^{ext}
=
\left(\frac{d\mathbf{H}_O}{dt}\right)_N
\]

Using the transport theorem:
\[
\left(\frac{d\mathbf{H}_O}{dt}\right)_N
=
\left(\frac{d\mathbf{H}_O}{dt}\right)_B
+
\bm{\omega}_{B/N}\times\mathbf{H}_O
\]

First compute the body-fixed derivative:
\begin{align}
\left(\frac{d\mathbf{H}_O}{dt}\right)_B
&=
I_1\ddot\theta\mathbf{b}_1
\\
&\quad
+
I_1(\ddot\psi\sin\theta+\dot\psi\dot\theta\cos\theta)\mathbf{b}_2
\\
&\quad
+
I_3(\ddot\phi+\ddot\psi\cos\theta
-\dot\psi\dot\theta\sin\theta)\mathbf{b}_3
\end{align}

Now compute the cross product term:
\[
\bm{\omega}_{B/N}\times\mathbf{H}_O
\]

Expanding term-by-term:
\begin{align}
\bm{\omega}\times\mathbf{H}
&=
(I_3-I_1)\dot\psi\sin\theta(\dot\phi+\dot\psi\cos\theta)\mathbf{b}_1
\\
&\quad
+
(I_1-I_3)\dot\theta(\dot\phi+\dot\psi\cos\theta)\mathbf{b}_2
\\
&\quad
+0\mathbf{b}_3
\end{align}

Therefore the Euler equations become:

\subsubsection*{$\mathbf{b}_1$ component}

\begin{align}
M_1
&=
I_1\ddot\theta
\\
&\quad
+
(I_3-I_1)
\dot\psi\sin\theta
(\dot\phi+\dot\psi\cos\theta)
\end{align}

\subsubsection*{$\mathbf{b}_2$ component}

\begin{align}
M_2
&=
I_1(\ddot\psi\sin\theta+\dot\psi\dot\theta\cos\theta)
\\
&\quad
+
(I_1-I_3)
\dot\theta(\dot\phi+\dot\psi\cos\theta)
\end{align}

\subsubsection*{$\mathbf{b}_3$ component}

\[
M_3
=
I_3(\ddot\phi+\ddot\psi\cos\theta
-\dot\psi\dot\theta\sin\theta)
\]

These are Euler's equations for the gyroscope.

\subsection*{(d) Steady Precession}

For steady precession:
\[
\dot\psi = \Omega_p = \text{constant}
\]
\[
\dot\phi = \Omega_s = \text{constant}
\]
\[
\theta = \text{constant}
\qquad\Rightarrow\qquad
\dot\theta = \ddot\theta = 0
\]

Substituting into the Euler equations:

\subsubsection*{$\mathbf{b}_1$ equation}

\[
M_1
=
(I_3-I_1)
\Omega_p\sin\theta
(\Omega_s+\Omega_p\cos\theta)
\]

\subsubsection*{$\mathbf{b}_2$ equation}

\[
M_2 = 0
\]

\subsubsection*{$\mathbf{b}_3$ equation}

\[
M_3 = 0
\]

Hence the external torque required to maintain steady precession is:
\[
\boxed{
M_1
=
(I_3-I_1)
\Omega_p\sin\theta
(\Omega_s+\Omega_p\cos\theta)
}
\]

\subsection*{(e) Lagrangian Formulation}

The kinetic energy of the rotor is:
\[
T
=
\frac12\bm{\omega}^T\mathbf{I}_O\bm{\omega}
\]

Substituting the angular velocity components:
\begin{align}
T
&=
\frac12 I_1\dot\theta^2
+
\frac12 I_1\dot\psi^2\sin^2\theta
\\
&\quad
+
\frac12 I_3(\dot\phi+\dot\psi\cos\theta)^2
\end{align}

Since the center of mass is fixed at $O$, the potential energy is constant:
\[
V = 0
\]

Thus the Lagrangian is:
\[
\mathcal{L}=T
\]

Applying Lagrange's equations:
\[
\frac{d}{dt}
\left(
\frac{\partial\mathcal{L}}{\partial\dot q_i}
\right)
-
\frac{\partial\mathcal{L}}{\partial q_i}
=0
\]

for:
\[
q_i = \{\psi,\theta,\phi\}
\]

produces the same equations obtained from the Newton-Euler formulation.

\subsection*{(f) Cyclic Coordinates}

The Lagrangian does not explicitly depend on:
\[
\psi,
\qquad
\phi
\]

Hence both coordinates are cyclic.

\subsubsection*{Momentum Associated with $\psi$}

\begin{align}
p_\psi
&=
\frac{\partial\mathcal{L}}{\partial\dot\psi}
\\
&=
I_1\dot\psi\sin^2\theta
+
I_3(\dot\phi+\dot\psi\cos\theta)\cos\theta
\end{align}

This is the angular momentum about the vertical axis.

\subsubsection*{Momentum Associated with $\phi$}

\begin{align}
p_\phi
&=
\frac{\partial\mathcal{L}}{\partial\dot\phi}
\\
&=
I_3(\dot\phi+\dot\psi\cos\theta)
\end{align}

This is the angular momentum along the rotor spin axis.

\subsection*{(g) Virtual Work Due to Gravity}

Suppose the center of mass is displaced from $O$ by:
\[
\mathbf{r}_{G/O}=\ell\mathbf{b}_3
\]

The gravitational force is:
\[
\mathbf{F}_g = -mg\mathbf{n}_3
\]

The virtual work is:
\[
\delta W
=
\mathbf{F}_g\cdot\delta\mathbf{r}_G
\]

The potential energy is:
\[
V = mg\ell\cos\theta
\]

The generalized forces are obtained from:
\[
Q_i = -\frac{\partial V}{\partial q_i}
\]

Thus:
\begin{align}
Q_\psi &= 0
\\
Q_\theta &= mg\ell\sin\theta
\\
Q_\phi &= 0
\end{align}